% !TeX root = ../main.tex

\section{Introduction}

Low Earth orbit (LEO) satellite networks are a promising architecture for
wide-area direct-access connectivity, but their downlink scheduling problem is
shaped by two coupled sources of variation. At the fast timescale, cells compete
for physical resource blocks (PRBs) and satellite transmit power under changing
demand. At the slow timescale, satellite visibility evolves with orbital motion,
which forces satellite-cell association decisions and may trigger frequent
handovers.

Existing LEO downlink allocation studies have emphasized fair resource
balancing across large service regions~\cite{leyva2023continent}, while
handover-aware formulations often use soft switching penalties
\cite{afif2025handover}. These approaches are valuable baselines, but two
practical gaps remain. First, PRB and power allocation should be optimized
jointly when evaluating fixed associations. Second, a soft handover penalty does
not guarantee that each cell stays within a deployment-level handover budget.
This motivates a two-timescale formulation that combines fast convex resource
allocation with slow hard-budget association optimization.

This paper focuses on the optimization core and leaves online learning under
uncertain demand as future work. The intended contributions are:
\begin{itemize}
\item a fast-timescale convex PRB--power allocation model for fixed LEO
associations;
\item a slow-timescale satellite-cell association formulation with hard
per-cell handover budgets;
\item exact and rolling-horizon solution methods that trade optimality for
runtime;
\item a reproducible simulation plan measuring throughput, proportional
fairness, worst-cell satisfaction, handover count, and solve time.
\end{itemize}
